\documentclass[12pt]{article}

% fonts

\usepackage[T1]{fontenc}
\usepackage{newtxtext}
\usepackage{bm}
\usepackage{amsthm} % before newtxmath b/c latex
\usepackage{newtxmath}
% \usepackage[bb=mt,cal=mt]{mathalpha}

\usepackage{amsmath}
\usepackage{amsfonts}
\usepackage{thmtools}
\usepackage{mathtools}

% spacing

\usepackage[parfill]{parskip}
\setlength{\parindent}{0pt}
\setlength{\parskip}{2ex plus 0.2ex minus 0.2ex}

% geometry of the page

\usepackage[
paper  = letterpaper,
left   = 1in,
right  = 1in,
top    = 1in,
bottom = 1in,
]{geometry}

% references

\usepackage{natbib}


\usepackage{tikzit}
\input{pm.tikzstyles}

\newcommand{\R}{\mathbb{R}}
\newcommand{\norm}[1]{\Vert#1\Vert}



\begin{document}

\begin{flushleft}
\textbf{%
    Approximation by Superpositions of a Sigmoidal Function (Cybenko, 1989)
} \\
Satvik Saha (ss7228) \\
\today
\end{flushleft}

\vspace{0.1in}


The flexibility of a neural network lies in its ability to approximate a very
large class of functions well, given an appropriate choice of activation
function $\sigma$; results of this nature are often termed as \emph{universal
approximation theorems}.
Here, we will only look at one-hidden-layer networks $\Omega\colon I \to \R$
with arbitrary width where $I := [0, 1]^n$.
These are maps of the form \[
    x \mapsto \sum_{j = 1}^N \alpha_j \sigma(y_j^\top x + \theta_j);
\] call this class of functions $\mathscr{S}_\sigma$.
\cite{cybenko89} shows that for continuous activation functions $\sigma$
satisfying $\lim_{t \to -\infty} \sigma(t) = 0$ and $\lim_{t \to \infty}
\sigma(t) = 1$, called \emph{sigmoidal} activation functions, the class
$\mathscr{S}_\sigma$ is dense in $(C(I), \norm{\cdot}_\infty)$.
This is also true of continuous $\sigma$ satisfying a more general
\emph{discriminatory} property: given a finite signed regular measure $\mu$ on
$I$, \[
    \int \sigma(y^\top x + \theta) \:d\mu(x) = 0\text{ for all } y \in \R^n, \theta \in \R
        \implies \mu = 0.
\] The proof of the result in \cite{cybenko89} is as follows: if $\mathscr{R}
:= \overline{\mathscr{S}_\sigma} \neq C(I)$, then one can produce a non-zero
bounded linear functional $L$ such that $L|_\mathscr{R} = 0$ via Hahn-Banach
extension.
Riesz Representation gives a finite signed regular measure $\mu_L$ such that $L
= \int \cdot \:d\mu_L$.
Consider the family of maps $\mathscr{H} := \{h_{y, \theta}: y \in \R^n, \theta
\in \R\}$ where each $h_{y, \theta}(x) = \sigma(y^\top x + \theta)$; since $
\mathscr{H} \subset \mathscr{S}_\sigma \subset \mathscr{R}$, we have
$L|_\mathscr{H} = 0$.
The discriminatory nature of $\sigma$ forces $\mu_L = 0$, contradicting $L \neq
0$.

For measurable functions $f$ on $I$ in general, given $\epsilon > 0$, Lusin's
theorem gives $f|_D = g$ continuous for some $D \subset I$ with Lebesgue
measure $m(D) > 1 - \epsilon$.
In this sense, neural networks can well approximate this broader class of
measurable functions.

This result does not hold if $\sigma$ is a polynomial.
Indeed, $h_{y, \theta}$ would also be a polynomial, and $\int h_{y, \theta}
\:d\mu = 0$ can only guarantee that finitely many moments of $\mu$ are zero,
hence $\sigma$ is not discriminatory.
More directly, $\mathscr{S}_\sigma$ would only contain polynomials of (at most)
a certain degree, and hence be a proper closed subspace of $C(I)$.

Fairly accessible classes of functions, such as polynomials and trigonometric
polynomials (finite linear combinations of $\cos(m\cdot), \sin(m\cdot)$), are
also dense in $C(I)$ via the Stone-Weierstrass theorem.
Recovering the right approximations would correspond to polynomial regression
and Fourier decomposition tasks.
What seems to make neural networks more attractive in practice is the
\emph{curse of dimensionality}; methods like polynomial regression are famously
sensitive when $n$ is large!
A good choice of activation function $\sigma$ can often make this regression
task more stable.


\vfill

\nocite{*}
\bibliographystyle{plainnat}

\begingroup
\renewcommand{\section}[2]{}%
\textbf{References}
\bibliography{week09}
\endgroup


\end{document}

%%% Local Variables:
%%% mode: latex
%%% TeX-master: t
%%% End:
